\SourceLeaf{167}your work known soon. You will surely earn gratitude for it, though hardly that
of the general public (which includes many a person regarded as
a skilled mathematician), for I am becoming more and more convinced
that the number of true geometers is extremely small and that most people
can neither judge the difficulties of work of this kind
nor even understand them --- but surely that of all
those whose judgment alone can be of real value to you. --- In
Braunschweig there is an émigré named \textsc{Chauvelot}, a
geometer of some ability, who claims to have established the theory
of parallel lines completely and will soon have his work printed,
but I expect nothing from him. In
\textsc{Hindenburg}'s Archiv, 9\textsuperscript{th} issue, there is
likewise a new attempt on the same subject by
a certain \textsc{Hauff}, which is beneath criticism. .\, .\, .\, .\, .\, .\, .

\medskip
\begin{center}
\rule{0.3\textwidth}{0.4pt}
\end{center}
\medskip

\begin{center}
\textsc{Gauss} to \textsc{Bolyai}. Braunschweig, 25 November 1804.
\end{center}

.\, .\, .\, .\, .\, .\, .\ I have read your essay with great interest and
attention, and took real pleasure in its genuine depth and insight.
But you do not want my empty praise, which might even
seem partial in a way simply because your
line of thought is very similar to the one by which I
formerly tried to untie this Gordian knot and have tried in vain
so far. You want only my honest, candid
judgment. It is that your procedure still does
\textls{not} 
satisfy me. I will try to bring the stumbling block I
still find in it (which belongs to the same group of reefs
on which my attempts have foundered so far) to light
as clearly as I can. I still hope
that those reefs will one day, and before my own
end, afford a passage. Meanwhile I have so
many other occupations on hand that I cannot give it
any thought at present. Believe me, I shall be truly
truly delighted if you get there before me and
